\subsection{Iterative Deepening Search (IDS)}

\textbf{Ý tưởng: }

Trong quá trình nghiên cứu các chiến lược tìm kiếm, hai thuật toán cơ bản là \textbf{Breadth-First Search (BFS)} và 
\textbf{Depth-First Search (DFS)} thể hiện rõ sự đánh đổi giữa \textit{tính tối ưu} và \textit{hiệu quả bộ nhớ} \cite{russell2021}:

\begin{itemize}
    \item \textbf{BFS} bảo đảm tìm được lời giải ngắn nhất (tối ưu) nhưng đòi hỏi phải lưu toàn bộ các nút trong hàng đợi ở cùng độ sâu. 
    Khi độ sâu của lời giải lớn, số lượng nút cần lưu trữ tăng theo cấp số nhân $O(b^d)$, dẫn đến \textbf{vấn đề tràn bộ nhớ}.
    \item \textbf{DFS} sử dụng bộ nhớ tuyến tính $O(bd)$ nhờ chỉ lưu các nút trên một đường đi hiện tại, 
    tuy nhiên nó có thể đi quá sâu trong không gian tìm kiếm và \textbf{bỏ qua lời giải nông hơn}, 
    hoặc thậm chí \textbf{rơi vào vòng lặp vô hạn}.
\end{itemize}

Để khắc phục nhược điểm này, \textbf{thuật toán Tìm kiếm Sâu dần Lặp (Iterative Deepening Search – IDS)} được phát triển 
dựa trên ý tưởng:
\begin{quote}
    “\textit{Duyệt không gian tìm kiếm theo chiều sâu giống DFS, nhưng theo thứ tự độ sâu tăng dần giống BFS.}”
\end{quote}

\textbf{Cấu trúc dữ liệu:} \texttt{stack}, biến giới hạn độ sâu \texttt{limit}, ánh xạ \texttt{visited[node] = parent}.

\textbf{Cơ chế thực hiện:}
\begin{itemize}
    \item IDS hoạt động bằng cách \textbf{gọi lại thuật toán DLS (Depth-Limited Search)} nhiều lần với giới hạn độ sâu $l$ tăng dần:
    \[
        l = 0, 1, 2, 3, \dots
    \]
    \item Trong mỗi lần lặp, IDS tiến hành duyệt theo chiều sâu nhưng chỉ tới mức giới hạn $l$. Các nút ở độ sâu $l$ được xem như không có con kế thừa.
    \item Sau khi hoàn thành một lượt tìm kiếm với giới hạn $l$, nếu mục tiêu chưa được tìm thấy, giới hạn được tăng lên $l+1$ và thuật toán khởi chạy lại từ nút gốc.
    \item Khi giới hạn $l$ đạt đến độ sâu $d$ của nút mục tiêu nông nhất, IDS sẽ dừng và trả về lời giải.
\end{itemize}

\textbf{Ghi chú triển khai: }
\begin{itemize}
    \item Trong khi DLS chỉ thực hiện một lần tìm kiếm với giới hạn cố định $l$, IDS tăng dần $l$ để đảm bảo tìm thấy lời giải mà không cần biết trước độ sâu.
\end{itemize}

\textbf{Đánh giá Hiệu suất:}
\begin{itemize}
    \item \textbf{Tính đầy đủ (Completeness):} Có, nếu hệ số nhánh $b$ là hữu hạn.
    \item \textbf{Tính tối ưu (Optimality):} Có, nếu chi phí bước là đồng nhất (mỗi bước đi có cùng chi phí). Khi đó, IDS tìm được lời giải ngắn nhất, tương tự BFS.
    \item \textbf{Độ phức tạp thời gian (Time Complexity):} $O(b^d)$, trong đó $d$ là độ sâu của lời giải nông nhất.  
    Mặc dù các nút ở các tầng trên được sinh ra nhiều lần, tổng số nút mở rộng vẫn tương đương BFS \cite{russell2021}:
    \[
        (d+1)b^0 + db^1 + (d-1)b^2 + \dots + 1b^d = O(b^d)
    \]
    \item \textbf{Độ phức tạp không gian (Space Complexity):} $O(bd)$, tuyến tính theo độ sâu, giống DFS.
    \item \textbf{Nhận xét:} 
    \begin{itemize}
        \item IDS được ưu tiên trong các bài toán có không gian tìm kiếm rất lớn hoặc khi độ sâu của lời giải chưa biết trước.
        \item DFS là trường hợp đặc biệt của DLS khi $l = \infty$, còn IDS chính là chuỗi các lần DLS tăng dần giới hạn.
    \end{itemize}
\end{itemize}